% !TEX program = xelatex
% 공통수학2 · Ⅰ. 도형의 방정식 · 2. 원의 방정식 · 중단원 마무리 (13차시)
% 학생용 활동지 (답안 없음)
\documentclass[11pt,a4paper]{article}
\usepackage{kotex}
\usepackage{iftex}
\ifPDFTeX\else
  \setmainhangulfont{Noto Serif CJK KR}
  \setsanshangulfont{Noto Sans CJK KR}
\fi
\usepackage[margin=2.1cm]{geometry}
\usepackage{parskip}
\usepackage{enumitem}
\usepackage{array}
\usepackage{tabularx}
\usepackage{xcolor}
\usepackage{amssymb}
\usepackage{tikz}
\usepackage[most]{tcolorbox}
\usepackage{fancyhdr}
\usepackage{titlesec}

\definecolor{goldc}{HTML}{B07C22}
\definecolor{goldstrong}{HTML}{8A5F16}
\definecolor{indigoc}{HTML}{2B3B57}
\definecolor{indigosoft}{HTML}{6E82A6}
\definecolor{linec}{HTML}{D6DACB}
\definecolor{surface2}{HTML}{E4E8DC}
\definecolor{writeline}{HTML}{C7CCBA}
\definecolor{inksoft}{HTML}{52604F}

\titleformat{\section}{\normalfont\sffamily\bfseries\large\color{indigoc}}{\thesection}{0.6em}{}
\titlespacing{\section}{0pt}{14pt}{6pt}
\newtcolorbox{goalbox}{enhanced, breakable, colback=white, colframe=goldc, boxrule=0.5pt, leftrule=3pt, arc=2pt, left=10pt, right=10pt, top=8pt, bottom=8pt}
\newtcolorbox{sheetbox}{enhanced, breakable, colback=white, colframe=linec, boxrule=0.6pt, arc=3pt, left=11pt, right=11pt, top=9pt, bottom=9pt}
\newcommand{\writespace}[1]{%
  \par\nobreak\vspace{3pt}
  \foreach \n in {1,...,#1} {%
    \noindent\textcolor{writeline}{\rule{\linewidth}{0.4pt}}\par\vspace{0.62cm}%
  }
}
\newcommand{\fillblank}[1]{\makebox[#1]{\hrulefill}}
\pagestyle{fancy}\fancyhf{}\renewcommand{\headrulewidth}{0pt}
\fancyfoot[L]{\footnotesize\color{inksoft} 공통수학2 · 2026학년도 2학기 · 지평선고등학교}
\fancyfoot[R]{\footnotesize\color{inksoft} 다음 시간 --- 3. 도형의 이동}

\begin{document}

\noindent
\begin{minipage}[b]{0.62\linewidth}
  {\sffamily\footnotesize\color{indigosoft} 공통수학2 · Ⅰ. 도형의 방정식}\\[2pt]
  {\sffamily\bfseries\Large 2. 원의 방정식 \textperiodcentered\ 중단원 마무리}
\end{minipage}\hfill
\begin{minipage}[b]{0.36\linewidth}
  \raggedleft\small\color{inksoft}
  반\ \fillblank{1.6cm} \quad 번호\ \fillblank{1.6cm}\\[4pt]
  이름\ \fillblank{3.6cm}
\end{minipage}

\vspace{4pt}
\noindent{\color{goldc}\rule{\linewidth}{2.2pt}}
\vspace{10pt}

\begin{goalbox}
\textbf{오늘의 목표}\\
9$\sim$12차시에서 배운 개념(원의 방정식, 원과 직선의 위치 관계, 접선)을 스스로 정리하고 점검한다.
\end{goalbox}

\section{개념 지도 채우기}

\begin{sheetbox}
\begin{tabularx}{\linewidth}{@{}>{\bfseries\small}p{4.8cm}X@{}}
\toprule
\textbf{개념} & \textbf{공식} \\
\midrule
원의 방정식(표준형) & \fillblank{6cm} \\[6pt]
원의 방정식(일반형이 원이 될 조건) & \fillblank{6cm} \\[6pt]
원과 직선 위치 관계 판단(판별식) & \fillblank{6cm} \\[6pt]
기울기가 주어진 접선 & \fillblank{6cm} \\[6pt]
원 위의 점에서의 접선 & \fillblank{6cm} \\
\bottomrule
\end{tabularx}
\end{sheetbox}

\section{기본 문제}

\begin{sheetbox}
\textbf{1.} 중심의 좌표가 $(1,2)$이고 반지름의 길이가 $\sqrt7$인 원의 방정식을 구하시오.
\writespace{2}

\textbf{2.} 원 $x^2+y^2+6x-8y+9=0$의 중심의 좌표 $(a,b)$와 반지름 $r$을 구하고, $a+b+r$의 값을 구하시오.
\writespace{2}

\textbf{3.} 원 $x^2+y^2=5$와 직선 $y=-3x+k$가 접하도록 하는 실수 $k$의 값을 구하시오.
\writespace{2}

\textbf{4.} 원 $x^2+y^2=17$ 위의 점 $(-4,1)$에서의 접선의 방정식을 구하시오.
\writespace{2}
\end{sheetbox}

\section{표준 \& 도전 문제}

\begin{sheetbox}
\textbf{5.} 방정식 $x^2+y^2-4x+2y+k=0$이 나타내는 도형이 반지름의 길이가 $\sqrt2$ 이상인 원이 되도록 하는 자연수 $k$의 개수를 구하시오.
\writespace{3}

\textbf{6. (서술형)} $x$축과 $y$축에 동시에 접하고 점 $(2,1)$을 지나는 원의 방정식을 모두 구하는 풀이 과정과 답을 쓰시오.
\writespace{4}

\textbf{7.} 원 $(x-2)^2+(y-5)^2=4$ 위의 점과 직선 $3x-4y-6=0$ 사이의 거리의 최댓값과 최솟값을 구하시오.
\writespace{3}
\end{sheetbox}

\section{마무리 성찰 --- 나의 성취 수준 점검}

\begin{sheetbox}
아래 항목 중 자신 있게 설명할 수 있는 것에 $\checkmark$ 표시해 보자.

\begin{itemize}[leftmargin=1.6em, itemsep=3pt]
  \item[$\square$] 중심과 반지름(또는 지름의 양 끝점)으로 원의 방정식을 구할 수 있다.
  \item[$\square$] 세 점을 지나는 원의 방정식을 구할 수 있다.
  \item[$\square$] 일반형이 원이 될 조건을 판단할 수 있다.
  \item[$\square$] 원과 직선의 위치 관계를 판단할 수 있다.
  \item[$\square$] 세 가지 상황에서 원의 접선의 방정식을 구할 수 있다.
\end{itemize}

\vspace{6pt}
{\footnotesize\color{inksoft} 가장 보충이 필요한 부분과 그 이유를 적어 보자.}
\writespace{2}
\end{sheetbox}

\end{document}
